% Drop-in LaTeX corrections for the TCSSOP manuscript.
% These blocks define the "fast schedule-aware scoring" step according to the
% Java implementation.

% ---------------------------------------------------------------------------
% 1) Replace vague references to "fast schedule-aware scoring function"
% ---------------------------------------------------------------------------

The fast schedule-aware scoring step is not an aggregate workload proxy and it
does not use a different dispatching rule from the final scheduler. It is a
lightweight execution mode of the same EDD-based constructive scheduler. For a
fixed chamber assignment \(X\) and sample vector \(S\), the scheduler constructs
the same ready-set sequence, applies the same chamber, humidity, voltage,
station, sample, and stage-state feasibility rules, and returns the resulting
total tardiness value. The computational saving comes from suppressing outputs
that are not needed during neighborhood screening: the fast mode does not build
the detailed operation-level schedule list or, when project-level lateness
information is not needed, the project-result list.

More specifically, a candidate move is scored as follows. First, chamber
instances are indexed by their assigned environment and by voltage capability.
Second, the job-based EDD scheduler repeatedly selects the best ready operation
from the project-level candidate pool. Third, each selected operation is placed
on the earliest feasible compatible station and sample. Finally, product
completion times and total tardiness are computed from the constructed state.
The returned score is therefore
\[
Score(X,S)=\sum_{p\in\mathcal{P}}\max(0,C_p(X,S)-d_p),
\]
where \(C_p(X,S)\) is obtained from the same deterministic scheduler used for
the final evaluation.

Two lightweight evaluation modes are used. The first mode returns only total
tardiness and is used when screening room-assignment moves, single-product
sample-size moves, and VNS shake candidates. The second mode returns total
tardiness and project-level completion/lateness values, but still skips the
operation-level schedule list; this mode is used when the Late--Early Rebalance
operator requires project-level lateness and earliness information. Accepted
outer-loop solutions are then re-evaluated in the schedule-building mode, which
constructs the full operation-level schedule and enables the independent
feasibility validation checks.

Thus, the term "fast" refers only to reduced output construction and, in some
calls, avoided defensive copying. It does not mean that the method evaluates a
candidate by an approximate capacity formula. All objective comparisons are
based on the same deterministic EDD schedule decoder, with the lexicographic
ordering \((T^{tot},S^{tot})\) applied after total tardiness is obtained.

% ---------------------------------------------------------------------------
% 2) Suggested replacement sentence for the local-search sections
% ---------------------------------------------------------------------------

% Suggested replacement sentence:
Candidate sample and chamber-assignment moves are screened by a lightweight
mode of the same EDD scheduler; this mode returns total tardiness without
building detailed schedule-output tables, and accepted outer-loop solutions are
redecoded with full schedule construction and validation.
